\renewcommand{\Anno}{2006}

\renewcommand{\NumeroProblema}{4}

\renewcommand{\Testo}{Si consideri il gioco seguente: su una
scacchiera $n\times n$ si mette una moneta nella casella in alto a sinistra e due giocatori $A$ e $B$ muovono a turno, cominciando da $A$, la moneta.  Ogni mossa consiste nello spostare la moneta di una casella, in orizzontale oppure in verticale, evitando di occupare le caselle già occupate in precedenza (sia da $A$ che da $B$).  Perde chi non riesce più a muovere la moneta in una casella ammissibile.

Determinare, in funzione del numero $n$, quale tra i 2 giocatori ha
una strategia vincente.}

\renewcommand{\Risposta}{Se $n$ è pari il giocatore $A$ ha una
strategia vincente, se $n$ è dispari il giocatore $B$ ha una
strategia vincente.}

\renewcommand{\Aiutino}{Supponiamo la scacchiera colorata nel modo
usuale e disposta in modo che la casella iniziale sia bianca.  Per $n$ pari è possibile costruire una corrisponenza biunivoca che ad ogni casella bianca associa una casella nera ad essa adiacente (con un lato in comune). In questo modo dopo ogni mossa di $B$ la pedina si trova su una casella nera ed $A$ può muovere nella bianca corrispondente.

Per $n$ dispari è possibile costruire una corrisponenza biunivoca che ad ogni casella nera associa una casella bianca (diversa da quella iniziale) ad essa adiacente.

In entrambi i casi la corrispondenza può essere ottenuta pensando
ad una tassellazione della scacchiera.}

\renewcommand{\SoluzioneUfficiale}{\TitoloSoluzione Dimostreremo che
per $n$ pari il giocatore $A$ ha una strategia vincente, mentre per $n$ dispari è il giocatore $B$ ad avere una strategia vincente.

Supponiamo intanto la scacchiera colorata nel modo usuale e disposta in modo che la casella iniziale sia bianca.  Il giocatore $A$ sposterà sempre la pedina da una casella bianca ad una casella nera, mentre il giocatore $B$ troverà sempre la pedina su una casella nera e la sposterà in una casella bianca.

Supponiamo ora che $n$ sia pari.  Allora è possibile ricoprire completamente la scacchiera usando delle mattonelle $2\times 1$, ciascuna delle quali occupa una casella bianca ed una casella nera. In questo caso il giocatore $A$ può, quando tocca a lui muovere, spostare la pedina dalla casella bianca in cui si trova alla casella nera che fa parte della stessa mattonella. Tale casella è sicuramente ammissibile perché essendo nera può essere raggiunta solo da $A$, e $A$ la raggiunge solo nel momento in cui $B$ decide di entrare in quella mattonella.

Supponiamo invece che $n$ sia dispari. Allora è possibile tassellare con mattonelle $2\times 1$ tutta la scacchiera meno la casella iniziale (basta ad esempio disporre le mattonelle in verticale nella prima colonna ed in orizzontale altrove). In questo caso i ruoli sono invertiti: $A$ sarà sempre il primo ad entrare in una casella nera di una certa mattonella, e $B$ potrà sempre replicare spostando la pedina nella casella bianca della mattonella stessa, fino alla vittoria.

Le seguenti figure rappresentano esempi di tassellazione per $n=6$ ed $n=7$.

\hfill
\psset{unit=5mm}
\pspicture(0,0)(6,6)
\psset{linewidth=0.6pt}
\newrgbcolor{scacchiera}{0.4 1 0.6}
	
\multips(0,0)(2,0){3}{\pspolygon*[linecolor=scacchiera](0,0)(1,0)(1,1)(0,1)(0,0)}
\multips(1,1)(2,0){3}{\pspolygon*[linecolor=scacchiera](0,0)(1,0)(1,1)(0,1)(0,0)}
\multips(0,2)(2,0){3}{\pspolygon*[linecolor=scacchiera](0,0)(1,0)(1,1)(0,1)(0,0)}
\multips(1,3)(2,0){3}{\pspolygon*[linecolor=scacchiera](0,0)(1,0)(1,1)(0,1)(0,0)}
\multips(0,4)(2,0){3}{\pspolygon*[linecolor=scacchiera](0,0)(1,0)(1,1)(0,1)(0,0)}
\multips(1,5)(2,0){3}{\pspolygon*[linecolor=scacchiera](0,0)(1,0)(1,1)(0,1)(0,0)}

\multido{\n=0+1}{7}{
\psline{-}(\n,0)(\n,6)
\psline{-}(0,\n)(6,\n)
}
\multips(0,0)(0,1){6}{\multips(0,0)(2,0){3}{\psframe[linewidth=1.5\pslinewidth](0.1,0.1)(1.9,0.9)}}
\endpspicture
\hfill
\pspicture(0,0)(7,8)
\psset{linewidth=0.6pt}
\newrgbcolor{scacchiera}{0.4 1 0.6}
	
\multips(1,0)(2,0){3}{\pspolygon*[linecolor=scacchiera](0,0)(1,0)(1,1)(0,1)(0,0)}
\multips(0,1)(2,0){4}{\pspolygon*[linecolor=scacchiera](0,0)(1,0)(1,1)(0,1)(0,0)}
\multips(1,2)(2,0){3}{\pspolygon*[linecolor=scacchiera](0,0)(1,0)(1,1)(0,1)(0,0)}
\multips(0,3)(2,0){4}{\pspolygon*[linecolor=scacchiera](0,0)(1,0)(1,1)(0,1)(0,0)}
\multips(1,4)(2,0){3}{\pspolygon*[linecolor=scacchiera](0,0)(1,0)(1,1)(0,1)(0,0)}
\multips(0,5)(2,0){4}{\pspolygon*[linecolor=scacchiera](0,0)(1,0)(1,1)(0,1)(0,0)}
\multips(1,6)(2,0){3}{\pspolygon*[linecolor=scacchiera](0,0)(1,0)(1,1)(0,1)(0,0)}

\multido{\n=0+1}{8}{
\psline{-}(\n,0)(\n,7)
\psline{-}(0,\n)(7,\n)
}
\multips(1,0)(0,1){7}{\multips(0,0)(2,0){3}{\psframe[linewidth=1.5\pslinewidth](0.1,0.1)(1.9,0.9)}
}

\multips(0,0)(0,2){3}{\psframe[linewidth=1.5\pslinewidth](0.1,0.1)(0.9,1.9)}
\endpspicture
\hfill\mbox{}
}

\renewcommand{\SoluzioniAlternative}{\TitoloSoluzioneAlternativaA }

\renewcommand{\Osservazioni}{\TitoloOsservazione }

\renewcommand{\Commenti}{
}